\documentclass[10pt]{article}
\usepackage[T1]{fontenc}
\usepackage[utf8]{inputenc}
\usepackage[french]{babel}
\usepackage{amsmath,amssymb,mathrsfs,mathpazo}
\usepackage{geometry}
\geometry{a4paper,margin=2cm}
\usepackage{graphicx,array,tabularx,tabularray,colortbl,multirow,multicol,booktabs,enumitem,pifont,ulem,fancybox,diagbox,cancel}
\usepackage{pstricks,pstricks-add,pst-plot,pst-tree,pst-node,pst-text,pst-func,pst-math,pst-bspline,pst-3dplot,pst-eucl,pst-solides3d}
\let\pscalculate\undefined
\usepackage{pst-bezier}
\usepackage{tikz,pgfplots,tkz-tab}
\usetikzlibrary{arrows,arrows.meta,patterns,calc,positioning,trees,decorations.markings,shapes.geometric,shapes.misc,shapes.arrows,intersections,decorations.pathmorphing,decorations.pathreplacing,backgrounds,shadows,babel}
\pgfplotsset{compat=1.18}
\usepgfplotslibrary{fillbetween}
\usepackage[np]{numprint}
\usepackage{listings,makecell,pas-tableur,twoopt}
\let\barre\undefined
\usepackage[pstricks]{tablvar}
\definecolor{darkgreen}{RGB}{0,100,0}
\providecommand{\red}{\color{red}}
\providecommand{\blue}{\color{blue}}
\providecommand{\green}{\color{green}}
\providecommand{\black}{\color{black}}
\providecommand{\white}{\color{white}}
\providecommand{\gray}{\color{gray}}
\providecommand{\cyan}{\color{cyan}}
\providecommand{\magenta}{\color{magenta}}
\providecommand{\yellow}{\color{yellow}}
\providecommand{\violet}{\color{violet}}
\providecommand{\orange}{\color{orange}}
\providecommand{\euro}{\texteuro}
\providecommand{\EUR}{\texteuro}
\providecommand{\ieme}{\textsuperscript{e}}
\providecommand{\degres}{\ensuremath{^\circ}}
\providecommand{\degree}{\ensuremath{^\circ}}
\providecommand{\vv}[1]{\overrightarrow{#1}}
\providecommand{\Coord}[1]{#1}
\providecommand{\decosix}{}
\providecommand{\starredbullet}{\ensuremath{\bullet}}
\pagestyle{empty}
\setlength{\parindent}{0pt}
\usepackage{adjustbox}\DeclareUnicodeCharacter{2236}{:}\DeclareUnicodeCharacter{2190}{\ensuremath{\leftarrow}}
\providecommand{\R}{\mathbb{R}}
\providecommand{\N}{\mathbb{N}}
\providecommand{\Z}{\mathbb{Z}}
\providecommand{\Q}{\mathbb{Q}}
\providecommand{\D}{\mathbb{D}}
\providecommand{\C}{\mathbb{C}}
\providecommand{\e}{{\mathrm{e}}}
\providecommand{\vect}[1]{\overrightarrow{#1}}
\providecommand{\vectt}[1]{\overrightarrow{\mathrm{#1}}}
\providecommand{\vvt}[1]{\overrightarrow{\mathrm{#1}}}
\providecommand{\np}[2][]{{#2~#1}}
\providecommand{\vv}[1]{{\overrightarrow{#1}}}
\renewcommand{\d}{{\mathrm{d}}}
\providecommand{\degres}{{^\circ}}
\providecommand{\degre}{{^\circ}}
\providecommand{\blue}{\color{blue}}
\providecommand{\red}{\color{red}}
\providecommand{\enskip}{\quad}
\providecommand{\llbracket}{\lbrack\!\lbrack}
\providecommand{\rrbracket}{\rbrack\!\rbrack}
\providecommand{\euro}{{\text{€}}}
\providecommand{\no}{{\text{n°}}}
\providecommand{\interval}[2]{{[#1;#2]}}
\providecommand{\intervFO}[2]{{[#1;#2[}}
\providecommand{\dotfill}{\ldots}

\providecommand{\ds}{\displaystyle}
\providecommand{\textding}[1]{#1}
\providecommand{\Oij}{\(\mathrm{O};\vec{\imath},\vec{\jmath}\)}
\providecommand{\Oijk}{\(\mathrm{O};\vec{\imath},\vec{\jmath},\vec{k}\)}
\providecommand{\Ouv}{\(\mathrm{O};\vec{u},\vec{v}\)}
\providecommand{\intervFF}[2]{\ensuremath{[#1;#2]}}
\renewcommand{\euro}{\ifmmode\text{\texteuro}\else\texteuro\fi}
\DeclareUnicodeCharacter{00B0}{\ifmmode^{\circ}\else\textdegree\fi}
\DeclareUnicodeCharacter{20AC}{\ifmmode\text{\texteuro}\else\texteuro\fi}
% Énoncés : Éducation nationale. Transcription et corrigés : APMEP, Denis Vergès et contributeurs. Découpage et classement : CoopMaths, Sébastien Lozano, Éric Elter et contributeurs. Adaptation de la présentation : Maths Explorer.
% CC BY-SA : https://coopmaths.fr/www/about/
% https://coopmaths.fr/alea/static/bac/2025/tex/bac_2025_06_sujet2_polynesie_2.tex
\begin{document}
\noindent\textit{APMEP / CoopMaths. CC BY-SA. Présentation adaptée par Maths Explorer.}\par\medskip
\medskip

On considère la fonction $f$ définie sur l'intervalle $]2~;~ +\infty[$ par 

\[f(x) = x \ln (x - 2).\]

Une partie de la courbe représentative $\mathcal{C}_f$ de la fonction $f$ est donnée ci-dessous.

\medskip

\begin{center}
\psset{xunit=1cm,yunit=0.2cm,arrowsize=2pt 3}
\begin{pspicture*}(-0.75,-17)(9,17)
\psaxes[linewidth=1.25pt,Dy=5,labelFontSize=\scriptstyle]{->}(0,0)(0,-17)(9,17)
\psplot[plotpoints=2000,linewidth=1.25pt,linecolor=blue]{2.0001}{9}{x 2 sub ln x mul}
\uput[ul](8,13){\blue $\mathcal{C}_f$}
\end{pspicture*}
\end{center}

\medskip

\begin{enumerate}
\item Conjecturer, à l'aide du graphique, le sens de variation de $f$ ses limites aux bornes
de son ensemble de définition ainsi que les éventuelles asymptotes.
\item Résoudre l'équation $f(x) = 0$ sur $]2~;~+\infty[$.
\item Calculer $\begin{gathered}\displaystyle\lim_{\substack{x \to 2\\x > 2}} f(x)\end{gathered}$.

Ce résultat confirme-t-il l'une des conjectures faites à la question 1. ?
\item Démontrer que pour tout $x$ appartenant à $]2~;~+\infty[$ :

\[f'(x) = \ln (x - 2) + \dfrac{x}{x - 2}.\]
\item On considère la fonction $g$ définie sur l'intervalle $]2~;~+\infty[$ par $g(x) = f'(x)$.
	\begin{enumerate}
		\item Démontrer que pour tout $x$ appartenant à $]2~;~+\infty[$, on a:

\[g'(x) = \dfrac{x - 4}{(x - 2)^2}.\]

		\item On admet que $\begin{gathered}\displaystyle\lim_{\substack{x \to 2\\x>2}} g(x) = + \infty\end{gathered}$ et que $\displaystyle\lim_{x \to + \infty} g(x) = +\infty$.

En déduire le tableau des variations de la fonction $g$ sur $]2~;~+\infty[$. On fera apparaître la valeur exacte de l'extremum de la fonction $g$.
		\item En déduire que, pour tout $x$ appartenant à $]2~;~+\infty[$,\: $g(x) > 0$.
		\item En déduire le sens de variation de la fonction $f$ sur $]2~;~+\infty[$.
	\end{enumerate}
\item Étudier la convexité de la fonction $f$ sur $]2~;~+\infty[$ et préciser les coordonnées d'un éventuel point d'inflexion de la courbe représentative de la fonction $f$.
\item Combien de valeurs de $x$ existe-t-il pour lesquelles la courbe représentative de $f$
admet une tangente de coefficient directeur égal à $3$ ?
\end{enumerate}

\bigskip

\end{document}
